\DocumentMetadata{
  lang=it-IT,
  pdfversion=2.0,
  pdfstandard=ua-2,
  tagging=on
}
\documentclass{article}
\usepackage[T1]{fontenc}
\usepackage[utf8]{inputenc}
\usepackage[italian]{babel}
\usepackage[accsupp]{axessibility}

\usepackage[
    unicode=true,
    hidelinks,
    pdfauthor={Università di Padova},
    pdftitle={},
    pdfsubject={Domande a risposta multipla},
    pdfkeywords={matematica, quiz, accessibile}
]{hyperref}

% Font senza grazie (TeX Gyre Heros = clone open source di Arial)
\usepackage{tgheros}
\renewcommand{\familydefault}{\sfdefault}

% Correzione legature problematiche per screen reader
\usepackage{microtype}
\DisableLigatures{encoding = *, family = *}

% Testo allineato a sinistra senza sillabazione
\usepackage{ragged2e}
\RaggedRight
\hyphenpenalty=10000
\exhyphenpenalty=10000

% Interlinea 1.5
\usepackage{setspace}
\setstretch{1.5}

% Dimensione titoli sezione a 18pt
\usepackage{titlesec}
\titleformat{\section}
    {\fontsize{18}{27}\selectfont\bfseries}
    {}{0em}{}

% Rimuove i numeri di pagina
\pagestyle{empty}

\begin{document}

\fontsize{18}{27}\selectfont

\section*{Domanda 1}

Luigi ha 2 figli di 15 e 11 anni.
Fra 18 anni la sua età sarà uguale alla somma delle età
che avranno i suoi figli.
Quanti anni ha oggi Luigi?

\begin{enumerate}
    \item[a)] 30
    \item[b)] Non si può dire
    \item[c)] 52
    \item[d)] 44
\end{enumerate}

\textbf{Risposta esatta: (d)}

\medskip

\section*{Domanda 2}

Il resto della divisione del polinomio
\[
    x^5 - 3x^4 + 3
\]
per $x + 1$ è:

\begin{enumerate}
    \item[a)] 1
    \item[b)] $-1$
    \item[c)] 3
    \item[d)] 0
\end{enumerate}

\textbf{Risposta esatta: (b)}

\medskip

\section*{Domanda 3}

Sia $T$ un triangolo rettangolo isoscele.
Allora la somma dei coseni degli angoli interni di $T$
è uguale a:

\begin{enumerate}
    \item[a)] 2
    \item[b)] 1
    \item[c)] $\sqrt{3}$
    \item[d)] $\sqrt{4}$
\end{enumerate}

\textbf{Risposta esatta: (d)}

\medskip

\section*{Domanda 4}

Il 30\% degli studenti iscritti a un corso universitario
ha superato l'esame relativo al primo appello.
Se, dei restanti studenti iscritti,
il 10\% supera l'esame al secondo appello,
gli studenti che devono ancora superare l'esame
dopo i primi due appelli saranno:

\begin{enumerate}
    \item[a)] Il 37\% del numero totale di studenti iscritti al corso
    \item[b)] Il 60\% del numero totale di studenti iscritti al corso
    \item[c)] Il 63\% del numero totale di studenti iscritti al corso
    \item[d)] Il 40\% del numero totale di studenti iscritti al corso
\end{enumerate}

\textbf{Risposta esatta: (c)}

\medskip

\section*{Domanda 5}

Trova il risultato della seguente disequazione:
\[
    \frac{2}{x} > 1
\]

\begin{enumerate}
    \item[a)] $x < 2$
    \item[b)] $0 < x < 2$
    \item[c)] $x > 2$
    \item[d)] $0 < x$ e $x > 2$
\end{enumerate}

\textbf{Risposta esatta: (b)}

\end{document}